% chapters/02/contoh.tex
\section{Dua Contoh Terapan}
\label{sec:contoh-bab-2}

Konsep-konsep pada bab ini akan lebih mudah dicerna apabila dilihat dalam
keadaan nyata. Dua contoh berikut memperlihatkan bagaimana konsep tersebut
bekerja pada satu bingkai permainan dan pada satu perhitungan sederhana.

\subsection{Contoh Pertama: Membaca Satu Bingkai Permainan}

\begin{contohbox}
Sebuah permainan dua dimensi berjalan pada 60 bingkai per detik. Pada awal
bingkai, seorang pemain berada pada posisi $(100, 300)$ dan bergerak ke kanan
dengan kecepatan 200 piksel per detik. Seekor musuh berada pada posisi
$(400, 300)$ dan berada pada keadaan mengejar. Sebuah koin berada pada posisi
$(160, 300)$.

\vspace{0.4em}
\textbf{Tahap 1: membaca masukan.} Program memeriksa peristiwa yang terjadi
sejak bingkai sebelumnya. Ternyata pemain menekan tombol panah kanan, sehingga
program mencatat bahwa arah gerak pemain adalah ke kanan. Penekanan tombol ini
tidak langsung memindahkan pemain; ia hanya mengubah keadaan yang akan dipakai
pada tahap berikutnya.

\vspace{0.4em}
\textbf{Tahap 2: memperbarui keadaan.} Program menghitung lama bingkai, yaitu
sekitar $0{,}0167$ detik. Dengan kecepatan 200 piksel per detik, pemain
berpindah sejauh $200 \times 0{,}0167 \approx 3{,}3$ piksel, sehingga posisinya
menjadi sekitar $(103{,}3, 300)$. Musuh menghitung arah menuju pemain, yaitu
ke kiri, lalu berpindah dengan cara yang sama. Setelah posisi keduanya
diperbarui, program memeriksa apakah persegi pembatas pemain beririsan dengan
persegi pembatas koin. Jarak keduanya masih sekitar 57 piksel, sehingga belum
bersentuhan dan skor belum bertambah.

\vspace{0.4em}
\textbf{Tahap 3: menayangkan.} Program menggambar latar, koin, musuh, dan
pemain pada posisi terbarunya, lalu menampilkan hasilnya di layar.

\vspace{0.4em}
Ketiga tahap itu berlangsung dalam waktu kurang dari 16,7 milidetik, kemudian
diulang dari awal. Skor, keadaan musuh, dan posisi seluruh
objek merupakan bagian dari satu kumpulan keadaan permainan yang dibawa dari
satu bingkai ke bingkai berikutnya.
\end{contohbox}

Contoh di atas memperlihatkan dua hal penting. Pertama, masukan tidak
mengubah posisi secara langsung, melainkan mengubah keadaan yang dipakai pada
perhitungan berikutnya. Kedua, urutan tahapannya tetap: masukan dibaca, keadaan
diperbarui, baru digambar. Menukar urutan tahapan itu menghasilkan perilaku
yang berbeda, misalnya gerak yang terasa tertinggal satu bingkai.

\subsection{Contoh Kedua: Menghitung Perpindahan pada Beberapa Laju Bingkai}

Bagian~\ref{sec:model-waktu} menyatakan bahwa kecepatan harus dinyatakan dalam
satuan per detik dan dikalikan dengan selang waktu. Contoh berikut menunjukkan
angka nyatanya.

\begin{table}[htbp]
\centering
\caption{Perpindahan objek berkecepatan 200 piksel per detik pada beberapa cara perhitungan}
\label{tab:contoh-dt}
{\setlength{\tabcolsep}{4pt}\footnotesize%
\begin{tabularx}{\textwidth}{@{}L{3.6cm} C{2.2cm} C{2.6cm} Y@{}}
\toprule
\headrow \textbf{Cara perhitungan} & \textbf{Lama bingkai} & \textbf{Perpindahan per bingkai} & \textbf{Jarak dalam satu detik} \\
\midrule
Selang waktu, 120 FPS & $8{,}3$ ms & $1{,}7$ piksel & 200 piksel \\
Selang waktu, 60 FPS & $16{,}7$ ms & $3{,}3$ piksel & 200 piksel \\
Selang waktu, 30 FPS & $33{,}3$ ms & $6{,}7$ piksel & 200 piksel \\
Tanpa selang waktu (ditambah tetap 5 piksel) & tidak dipakai &
5 piksel & 150 piksel pada 30 FPS; 300 piksel pada 60 FPS \\
\bottomrule
\end{tabularx}}
\end{table}

Tiga baris pertama pada Tabel~\ref{tab:contoh-dt} memperlihatkan sifat yang
diinginkan: perpindahan pada setiap bingkai berbeda, tetapi jarak yang ditempuh
dalam satu detik selalu sama, yaitu 200 piksel. Baris terakhir
memperlihatkan akibat apabila selang waktu diabaikan. Pada 30 FPS objek
bergerak sejauh 150 piksel dalam satu detik, sedangkan pada 60 FPS ia bergerak
sejauh 300 piksel, dua kali lebih cepat. Permainan menjadi jauh lebih sulit
pada komputer yang lebih cepat, dan perilakunya berubah setiap kali perkakas
atau beban kerjanya berubah.

Keadaan sebaliknya juga perlu diperhatikan. Apabila kecepatan dinyatakan
dalam satuan per bingkai, penambahan beban kerja pada tahap penyajian akan
memperlambat laju bingkai, dan seluruh permainan ikut melambat. Pemain melihat
gerak yang tersendat, dan menganggap programnya berat, padahal yang terjadi
adalah perhitungan yang mengaitkan gerak pada bingkai, bukan pada waktu.
